% !TEX root = userguide.tex

\chapter{Solvers}
\pgst{tfemch}

\def\chaptertitle{Solvers}
\def\headdate{18th October 2023}

As we already mentioned in the introduction chapter: \tfem\ does not have
a solver included.
In this chapter we give some possibilities for using external solvers with
\tfem.

\section{Sparse direct solvers from the HSL library: MA41 and MA57}
\label{sec:hslsparse}

For symmetric and unsymmetric matrices we can use the MA57 and MA41
sparse direct solver from the HSL library\footnote{
The HSL library must be licensed \cite{HSL}. Academics can obtain a personal license for free.
See \href{http://www.hsl.rl.ac.uk/}{http://www.hsl.rl.ac.uk/} for more information.}, respectively (see Sec.~\ref{sec:hsl_addon}).
The HSL library has been written in Fortran, which makes interfacing with \tfem\ easy. In this way, it is possible to use the powerful \texttt{-C=undefined} checking option of the nagfor compiler for \tfem.

For the Poisson and Stokes problems of Chapter~\ref{ch:examples},
the system matrix is symmetric.
For such problems it is both
CPU and memory efficient to store only the upper-triangle of the matrix and
use the symmetric sparse direct solver MA57. This can be achieved as follows:
\begin{verbatim}
  ....
  use hsl_ma57_m
  ...
  call create_sysmatrix_structure ( sysmatrix, mesh, problem, symmetric=.true. )
  ...
  call solve_system_ma57 ( sysmatrix, rhsd, sol )
  ...
\end{verbatim}
The \texttt{symmetric=.true.} keyword tells the routine to store only the
upper-triangular part of the matrix. The same keyword is available for the
lower-level routine \texttt{create\_sysmatrix\_structure\_base}. For an
unsymmetric matrix we should specify \texttt{symmetric=.false.} or
omit the keyword \texttt{symmetric} altogether:
\begin{verbatim}
  ....
  use hsl_ma41_m
  ...
  call create_sysmatrix_structure ( sysmatrix, mesh, problem )
  ...
  call solve_system_ma41 ( sysmatrix, rhsd, sol )
  ...
\end{verbatim}
If you want to use both MA41 and MA57 in the same program unit, you can `use'
both modules like:
\begin{verbatim}
  ....
  use hsl_ma41_m
  use hsl_ma57_m
  ...
\end{verbatim}
Both MA41 and MA57 can scale the matrix\footnote{For MA57 you need a recent version. The HSL2002 version does not have scaling.}, which can make solving systems with large differences in element size and/or fluid properties much faster.

It is strongly advised to combine both MA41 and MA57 with \metis\
renumbering
(see Sec.~\ref{sec:metis5}) to change the pivot order to reduce the fill-in, especially for larger problems with Lagrange multipliers, additional unknowns and/or connections. The speedup and memory reduction can be significant. An example for MA41 is given in Sec.~\ref{sec:poisson24}. For MA57 you need a recent version that automatically uses \metis\ if linked in. The usual MA57 from HSL2002 as supplied to students cannot be used together with \metis\ in \tfem.

\section{Other sparse direct solvers}

Interfaces to other sparse direct solvers are available as an add-on:
\begin{itemize}
   \item MUMPS, see Sec.~\ref{sec:mumps}.
         Sequential and parallel version using MPI.
         Public domain.
   \item PARDISO, see Sec.~\ref{sec:pardiso}. 
         Sequential and parallel version using
         OPENMP. Free for academic use. Also Intel's MKL contains PARDISO.
   \item UMFPACK, see Sec.~\ref{sec:umfpack}. UMFPACK is free software and written in C.
   \end{itemize}
UMFPACK assumes the matrix is unsymmetric\footnote{A symmetric matrix must be stored as if it is unsymmetric.}, whereas MUMPS and PARDISO can be used for both symmetric and unsymmetric matrices. Note, that when the symmetric solver is used, the structure of
\texttt{sysmatrix} must be created using \texttt{symmetric=.true.} as explained
in Sec.~\ref{sec:hslsparse} for MA57.

\section{Iterative solvers}

Interfaces to iterative solvers are available as an add-on:
\begin{itemize}
   \item SPARSKIT2, see Sec.~\ref{sec:sksolve}. Public domain. 
   \item MKL, see Sec.~\ref{sec:mklgmres}. Free to use.
   \item HSL, see Sec.~\ref{sec:hsl_extra_addon}. Free for academics.
    Includes algebraic multigrid (AMG) preconditioning.
\end{itemize}

%In 2D the direct solvers are usually faster than iterative solvers, provided
%enough memory is available.
In 3D the full LU decomposition requires a huge
amount of both memory and CPU time. Therefore, using iterative solvers is
the only option in many cases.
If algebraic multigrid (AMG) preconditioning is applicable, $\mathcal{O}(N)$
solvers can be achieved.

\section{Eigenvalue solvers}

Interfaces to eigenvalue solvers are available as an add-on:
\begin{itemize}
   \item FEAST, see Sec.~\ref{sec:feast}. Public domain. 
\end{itemize}

